\subsection{Simulation of a complete beamline}

Designing and optimizing these beamlines is a complex task, because every optical element introduces aberrations, scattering, and intensity losses that interact in non-trivial ways. This is where ray tracing simulation of the full beamline, including the particularities of every optical element becomes indispensable.

The core challenge is that real synchrotron beamlines are expensive, time-consuming to build, and often inaccessible during early design phases. Ray tracing allows scientists and engineers to:

\begin{itemize}
    \item Predict beam properties (spot size, divergence, flux, energy resolution) at any point along the beamline before a single component is installed.
    \item Optimize optical layouts by virtually testing hundreds of mirror geometries, monochromator configurations, or aperture settings.
    \item Study the tolerances in the alignment of each element, and the optical effect of the misalignment. 
    \item Anticipate and understand the possible problems and risks.
    \item Train new users by providing a realistic virtual environment.
\end{itemize}


Diagnose aberrations — understanding why a beam is elongated, defocused, or contaminated with higher harmonics.
